% sections/results.tex
% All internal phase labels, model filenames, and run IDs removed from main text.
% Experiments referred to by what they test, not by internal planning labels.

\subsection{The GCSF/EPO Ratio, Not EPO Dose Alone, Controls Erythroid
Commitment Probability}
\label{sec:ratio_sweep}

The central quantitative result of this work is that erythroid commitment
probability is governed by the \emph{ratio} of GCSF to EPO receptor occupancy,
not by the absolute EPO concentration. To characterise the bifurcation
structure, we ran a ten-condition GCSF sweep at fixed EPO with $N=50$ stochastic
replicates per condition. The result is an abrupt, all-or-nothing switch
(Table~\ref{tab:commitment_D}; Fig.~\ref{fig:p_erythroid}): at all GCSF doses
below the threshold, erythroid commitment fraction exceeds 0.92; at doses above
the threshold, the fraction collapses to zero. The commitment threshold lies
between 2.2 and 2.5~mM GCSF at EPO\,=\,1.0~mM (ratio 2.2--2.5), with no
intermediate values in 500 replicates.

\begin{table}[ht]
\caption{GCSF/EPO ratio sweep ($N=50$ per condition, EPO\,=\,1.0~mM fixed,
$T=7200$~s). GATA1$_\text{nuc}$ and PU.1$_\text{nuc}$ are mean\,$\pm$\,SD across
50 replicates at $t=3600$~s. P(ERY) is the fraction of replicates with
GATA1\_nuc\,$>$\,1.5\,$\times$\,PU.1\_nuc; 95\% CI uses the Wilson score. The
transition between complete erythroid and complete myeloid commitment occurs
between GCSF\,=\,2.2 and GCSF\,=\,2.5~mM (GCSF/EPO\,$\in$\,(2.2,\,2.5)).}
\label{tab:commitment_D}
\centering
\small
\resizebox{\linewidth}{!}{\begin{tabular}{@{}cccccc@{}}
\toprule
\textbf{GCSF (mM)} & \textbf{GCSF/EPO} & \textbf{GATA1\textsubscript{nuc} (mM)}
& \textbf{PU.1\textsubscript{nuc} (mM)} & \textbf{P(ERY)} & \textbf{95\% CI} \\
\midrule
0.5  & 0.5 & $9.36 \pm 1.81$ & $0.003 \pm 0.018$ & 1.00 & [0.929, 1.00] \\
1.0  & 1.0 & $9.25 \pm 1.43$ & $0.004 \pm 0.023$ & 1.00 & [0.929, 1.00] \\
1.2  & 1.2 & $9.14 \pm 1.97$ & $0.133 \pm 0.927$ & 0.98 & [0.895, 0.996] \\
1.35 & 1.35 & $8.75 \pm 2.89$ & $0.710 \pm 2.43$ & 0.92 & [0.812, 0.968] \\
1.5  & 1.5 & $9.31 \pm 2.46$ & $0.351 \pm 1.71$ & 0.96 & [0.865, 0.989] \\
2.0  & 2.0 & $9.24 \pm 2.10$ & $0.170 \pm 1.16$ & 0.98 & [0.895, 0.996] \\
2.2  & 2.2 & $8.81 \pm 2.26$ & $0.403 \pm 1.96$ & 0.96 & [0.865, 0.989] \\
2.5  & 2.5 & $0.005 \pm 0.016$ & $8.83 \pm 1.20$ & 0.00 & [0.000, 0.071] \\
2.8  & 2.8 & $0.003 \pm 0.013$ & $8.81 \pm 1.42$ & 0.00 & [0.000, 0.071] \\
3.0  & 3.0 & $0.585 \pm 2.29$ & $8.30 \pm 2.54$ & 0.06 & [0.021, 0.162] \\
\bottomrule
\end{tabular}}
\end{table}

At GCSF\,=\,1.35~mM the commitment variance is maximal (GATA1$_\text{nuc}$ SD
= 2.89; PU.1$_\text{nuc}$ SD = 2.43) with 46/50 erythroid replicates---a
genuine mixed-fate competition zone in which individual cell outcomes are
stochastic. The erythroid attractor (GATA1$_\text{nuc}\approx 9$~mM, PU.1$_\text{nuc}\approx 0$) and myeloid attractor (GATA1$_\text{nuc}\approx 0.005$~mM,
PU.1$_\text{nuc}\approx 8.8$~mM at GCSF\,=\,2.5) are nearly mirror-symmetric
(Fig.~\ref{fig:waddington}), consistent with a genuine bistable switch whose
basin boundary lies at GCSF/EPO\,$\approx$\,2.3.

\begin{figure}[!ht]
  \centering
  \includegraphics[width=0.75\linewidth]{fig2_p_erythroid}
  \caption{\textbf{P(erythroid) vs GCSF/EPO ratio.} Erythroid commitment
  fraction at $t=3600$~s across ten GCSF conditions at EPO\,=\,1.0~mM fixed
  ($N=50$ per condition). P(ERY) is high and roughly flat for GCSF/EPO\,$\leq$\,2.2,
  then switches abruptly to zero at GCSF/EPO\,=\,2.5.}
  \label{fig:p_erythroid}
\end{figure}

The sharpness of this switching curve reflects two deep, well-separated potential wells rather than a shallow continuum of states: the erythroid attractor (GATA1$_\text{nuc}\approx 9$~mM, PU.1$_\text{nuc}\approx 0$) and myeloid attractor (PU.1$_\text{nuc}\approx 8.8$~mM, GATA1$_\text{nuc}\approx 0$) are nearly mirror-symmetric in depth, with a sharp potential barrier whose projection onto the GCSF/EPO axis aligns with the commitment midpoint at ratio\ $\approx$\,2.3.

\begin{figure}[!ht]
  \centering
  \includegraphics[width=0.75\linewidth]{fig1_waddington}
  \caption{\textbf{Erythroid and myeloid basins of attraction.} Approximate
  Waddington potential on the GATA1\textsubscript{nuc}$\times$PU.1\textsubscript{nuc}
  plane; attractor wells estimated from the maximum-variance condition
  (GCSF\,=\,1.35~mM, $N=50$).}
  \label{fig:waddington}
\end{figure}

The EPO-only dose-response result reported in a prior sweep (Section~\ref{sec:epo_dose})
is a direct consequence of this ratio-control principle: varying EPO in a narrow
window that barely changes the GCSF/EPO ratio produces a system that appears
EPO-insensitive. The apparent insensitivity is a property of the measurement
window, not of the circuit.

\subsection{EPO Dose Alone Does Not Determine Fate at Single-Cell Copy Numbers}
\label{sec:epo_dose}

To characterise the system in the single-cell regime, we ran $N=50$ stochastic
replicates at each of eight EPO concentrations spanning a 7\% range at fixed
GCSF, operating at one-hundredth of standard initial conditions to place the
model in the physiological single-cell copy-number regime. Erythroid commitment
fraction was statistically flat across the entire EPO range (60--70\%;
Kruskal--Wallis $p>0.05$), with high replicate-to-replicate variance in both
GATA1\textsubscript{nuc} and PU.1\textsubscript{nuc} at the end of each
simulation. No dose-dependent increase in commitment fraction was detectable.

This result is a direct consequence of the GCSF/EPO ratio control demonstrated
above: within a 7\% EPO window the ratio shifts by only \textasciitilde7\%, which
is far insufficient to cross the bifurcation boundary. Any experiment that fixes
one cytokine and sweeps the other in a narrow neighbourhood of the system's
operating point will appear cytokine-insensitive---the informative quantity is
always the ratio.

The energy nucleotide species provided an internal consistency control for all
single-cell simulations: total adenine nucleotide (ATP + ADP) and total guanine
nucleotide (GTP + GDP) were conserved to machine precision across all 400
replicates, confirming numerical integrity throughout.

\subsection{A Three-Layer Cascade Architecture}
\label{sec:layers}

Analysis of the mean trajectory---the average over all 50 replicates in each
condition---reveals that the dynamics from stimulation to settled state are
organised into three functionally distinct layers whose timing and EPO-sensitivity
differ sharply from one another (Fig.~\ref{fig:cascade}, Table~\ref{tab:layers}).

\begin{figure}[!ht]
  \centering
  \includegraphics[width=0.85\linewidth]{fig7_preemption_cascade}
  \caption{\textbf{Signal hierarchy preemption cascade.} Four panels show
  sequential activation---EPO forcing, EPOR binding, pGATA1 commitment pulse,
  and GATA1/PU.1 execution divergence---over the first 400~s of stimulation.}
  \label{fig:cascade}
\end{figure}

\subsubsection*{Layer 1: Reception}

The reception layer comprises cytokine receptor binding, internalisation, and
recycling at the cell surface. Its defining characteristic is a structural
erythroid bias that is essentially invariant across the tested EPO range: EPO
receptors are roughly four times more occupied than GCSF receptors regardless of
EPO concentration, because GCSF is already saturating its receptor at the fixed
dose used throughout. Doubling EPO concentration changes EPOR occupancy by only
a few percent---an effect too small to meaningfully shift downstream signalling.
The reception layer sets a coarse, constitutive erythroid bias and passes it
downstream unchanged.

\subsubsection*{Layer 2: Decision}

The decision layer is where GATA1 and PU.1 compete for transcriptional
dominance. Both genes fire at near-maximal rate from the first instant of
stimulation---the system is transcriptionally primed and requires no activation
lag. A persistent 4--5:1 GATA1-to-PU.1 transcription advantage provides the
mean erythroid bias, but because the absolute number of mRNA molecules involved
is very small, a single stochastic fluctuation---one extra PU.1 mRNA
transcript---can override this average and tip the cell toward the myeloid fate.
EPO dose does not modulate the transcriptional asymmetry; it is structurally set.
The decision layer is therefore the locus of stochastic fate specification
(Fig.~\ref{fig:divergence}).

\begin{figure}[!ht]
  \centering
  \includegraphics[width=0.85\linewidth]{fig3_divergence}
  \caption{\textbf{Stochastic divergence (mean $\pm$1\,SD envelopes).}
  GATA1\textsubscript{nuc} and PU.1\textsubscript{nuc} mean\,$\pm$\,1\,SD over
  $N=50$ replicates, $t=0$--3600~s; wide envelopes reflect genuine bimodality
  driven by decision-layer noise.}
  \label{fig:divergence}
\end{figure}

\subsubsection*{Layer 3: Execution}

Once GATA1 accumulates in the nucleus, the execution layer fires: GATA1 is
phosphorylated by a receptor-coupled kinase, and the phosphorylated form drives
ubiquitin-mediated degradation of PU.1. This cascade activates within the first
minute and completes within roughly two minutes. Critically, it is
EPO-dose-insensitive: once triggered, it runs to the same endpoint at the same
speed regardless of how much EPO is present. The execution layer converts an
upstream binary decision into a cytokine-maintained chemical record.

\begin{table}[ht]
\caption{Three-layer cascade summary. Functional properties of each layer.
GATA1:PU.1 transcription ratio refers to mean cumulative firing counts;
activation and peak times are from the mean trajectory.}
\label{tab:layers}
\centering
\small
\begin{tabularx}{\linewidth}{@{}lXXX@{}}
\toprule
\textbf{Layer} & \textbf{Main components} & \textbf{EPO-dose sensitivity} &
\textbf{Key property} \\
\midrule
Reception & Cytokine receptor binding, internalisation, recycling &
Very low ($<$6\% change across 7\% EPO range) & Structural 4:1 EPOR/GCSFR
bias; set by receptor abundance \\
\addlinespace
Decision & GATA1 \& PU.1 transcription, mRNA export, translation,
nuclear import & None (ratio 4--5:1 across all EPO doses) & Stochastic noise
at single-digit mRNA counts determines fate \\
\addlinespace
Execution & GATA1 phosphorylation, PU.1 ubiquitin degradation & None &
Activates $\sim$48~s; peaks $\sim$77~s; EPO-maintained \\
\bottomrule
\end{tabularx}
\end{table}

The net result on the mean trajectory (Fig.~\ref{fig:mean_trajectory}):
GATA1\textsubscript{nuc} rises continuously while PU.1\textsubscript{nuc} erodes
over 3600~s, a signature that is EPO-dose-insensitive across all tested
conditions.

\begin{figure}[!htb]
  \centering
  \includegraphics[width=0.85\linewidth]{fig4_mean_trajectory}
  \caption{\textbf{Mean trajectory: GATA1 rise, PU.1 erosion.}
  Population-mean GATA1\textsubscript{nuc} and PU.1\textsubscript{nuc} over
  $t=0$--3600~s ($N=50$); the continuous GATA1 rise and PU.1 decline are the
  cumulative output of all three cascade layers.}
  \label{fig:mean_trajectory}
\end{figure}

\FloatBarrier
\subsection{Cytokine-Dependent Priming Crossover}
\label{sec:priming}

From the mean trajectory, we identified the point at which phosphorylated
nuclear GATA1 first exceeds the nuclear PU.1 protein level. This crossing
occurs at approximately 146 seconds after EPO stimulation (Fig.~\ref{fig:ratchet},
Table~\ref{tab:ratchet}) and marks an important priming event: from this moment
the EPOR-driven phosphorylation cascade can sustain PU.1 degradation faster than
PU.1 is replenished, creating a self-reinforcing positive feedback as long as EPO
receptor signalling persists. However, as demonstrated by the cytokine-withdrawal
experiment (Section~\ref{sec:withdrawal}), this state is \emph{not} an irreversible
molecular record: it collapses if the signal is removed.

\begin{figure}[!ht]
  \centering
  \includegraphics[width=0.85\linewidth]{fig5_ratchet}
  \caption{\textbf{Execution-layer priming onset ($t=0$--400\,s).}
  pGATA1\textsubscript{nuc} and PU.1\textsubscript{nuc} mean trajectories in
  the first 400~s; the crossover at $t\approx146$~s marks the cytokine-dependent
  priming point.}
  \label{fig:ratchet}
\end{figure}

The five functional milestones that define this kinetic window are tabulated in Table~\ref{tab:ratchet}, providing quantitative anchors for the cascade sequence: the time of execution-layer activation ($\sim$48~s), ubiquitin-cascade peak ($\sim$77~s), priming crossover ($\sim$146~s), and the final committed steady state at 3600~s.

\begin{table}[ht]
\caption{Commitment state at key timepoints (mean trajectory, continuous
stimulation). GATA1$_\text{nuc}$ includes total unphosphorylated plus
phosphorylated nuclear protein; phosphorylated fraction shown separately.}
\label{tab:ratchet}
\centering
\small
\resizebox{\linewidth}{!}{\begin{tabular}{@{}lcccc@{}}
\toprule
\textbf{Event} & \textbf{Time} & \textbf{GATA1\textsubscript{nuc}} &
\textbf{pGATA1 fraction} & \textbf{PU.1\textsubscript{nuc}} \\
\midrule
EPO stimulation & 0~s & baseline & --- & \textasciitilde\,0.416 \\
Execution layer activated & $\sim$48~s & rising & --- & declining \\
PU.1 ubiquitin cascade peaks & $\sim$77~s & rising & --- & declining \\
GATA1/PU.1 priming crossover & $\sim$146~s & \multicolumn{2}{c}{pGATA1 = PU.1}
& declining \\
Committed steady state & 3600~s & 1.04 & 59\% & 0.26 \\
\bottomrule
\end{tabular}}
\end{table}

To assess robustness of the 146~s prediction to uncertainty in the
phosphorylation rate constant, a local sensitivity analysis rescaled the relevant
parameter by factors $\alpha \in \{0.50, 0.80, 1.20, 1.50\}$. A $\pm$20\%
perturbation shifts the priming-crossover time by only $\pm$3\%, while a $\pm$50\%
perturbation shifts it by $-40\%$ to $+49\%$ (87--217~s). The qualitative
prediction---an early cytokine-dependent priming event within the first two
minutes of EPOR activation---is preserved across the entire range tested
(Fig.~\ref{fig:ratchet_sensitivity}).

\begin{figure}[!htb]
  \centering
  \includegraphics[width=0.82\linewidth]{fig9_ratchet_sensitivity}
  \caption{\textbf{Priming-crossover time robustness.}
  \textbf{A,} pGATA1$_{\text{nuc}}$ and PU.1$_{\text{nuc}}$ trajectories for
  $t = 0$--400~s under five phosphorylation-rate scale factors; vertical lines
  mark the corresponding crossover times.
  \textbf{B,} Priming-crossover time vs scale factor; dashed line marks the
  baseline (146~s). Even at $\pm50\%$ perturbation the priming crossover occurs
  within 4 minutes of EPO stimulation.}
  \label{fig:ratchet_sensitivity}
\end{figure}

\subsection{EPO Withdrawal Reverses Erythroid Commitment}
\label{sec:withdrawal}

To test whether the erythroid attractor represents a true irreversible state, we
performed a cytokine-withdrawal experiment. EPO was applied at 1.0~mM for 500~s
(well past the $t\approx146$~s priming crossover), then removed by a timed
environment event. The trajectory was recorded to $t=3600$~s
(Table~\ref{tab:withdrawal}).

Following EPO withdrawal the erythroid state was not maintained. By $t=1000$~s
the GATA1/PU.1 ratio had fallen from 5.97 to 3.00; by $t=2000$~s
GATA1$_\text{nuc}\approx 0$ and PU.1$_\text{nuc}=1.996$, constituting a complete
fate reversal to the myeloid attractor. The final state at $t=3600$~s was
GATA1$_\text{nuc}\approx 0$, PU.1$_\text{nuc}=1.653$.

\begin{table}[ht]
\caption{Cytokine-withdrawal trajectory: EPO\,=\,1.0~mM for $t=0$--500~s, then
EPO\,=\,0. GATA1 dominance established at $t=490$~s collapses progressively
following cytokine withdrawal.}
\label{tab:withdrawal}
\centering
\small
\begin{tabularx}{\linewidth}{@{}lccc@{}}
\toprule
\textbf{Event} & \textbf{Time} & \textbf{GATA1\textsubscript{nuc}} &
\textbf{PU.1\textsubscript{nuc}} \\
\midrule
EPO stimulation      & 0~s      & baseline & \textasciitilde\,0.385 \\
Peak erythroid (pre-withdrawal) & 490~s & 10.91 & \textasciitilde\,0 \\
EPO withdrawn        & 500~s    & ---      & --- \\
GATA1/PU.1 ratio = 3.00 & 1000~s & declining & 1.33 \\
Full fate reversal   & $\sim$2000~s & $\approx$0 & 1.996 \\
Final myeloid steady state & 3600~s & $\approx$0 & 1.653 \\
\bottomrule
\end{tabularx}
\end{table}

The mechanistic basis is direct: EPOR occupancy is the dominant multiplicative
driver of GATA1 transcription. Without EPO, EPOR occupancy decays to zero,
reducing GATA1 transcription to its bare self-activation rate, which is
insufficient to replenish GATA1 against constitutive degradation. PU.1 then
rebounds via unimpeded basal transcription.

This result demonstrates that the erythroid attractor in this model is an
EPOR-signalling-dependent state, not an epigenetically sealed absorbing state.
True transcriptional irreversibility requires chromatin remodelling that deposits
a memory of commitment independent of ongoing receptor occupancy---a layer absent
from the present model.

\subsection{A Narrow, Non-Monotone EPO Pulse Window Determines Post-Withdrawal
Fate}
\label{sec:pulse}

To map the relationship between EPO pulse duration and post-withdrawal fate
maintenance, we performed a systematic pulse-duration sweep. EPO was applied at
1.0~mM and removed at times spanning 60 to 1800~s; each trajectory was recorded
to $t=3600$~s (Table~\ref{tab:pulse}).

\begin{table}[ht]
\caption{Pulse-duration sweep. EPO\,=\,1.0~mM applied for $t_{\mathrm{pulse}}$~s
then removed; GATA1$_\text{nuc}$ and PU.1$_\text{nuc}$ at $t=3600$~s.
The fate outcome is non-monotone: only the 120~s and 150~s pulses maintain
erythroid fate. Ghost-attractor reversal time scales systematically with pulse
length.}
\label{tab:pulse}
\centering
\small
\resizebox{\linewidth}{!}{\begin{tabular}{@{}cccclc@{}}
\toprule
\textbf{$t_{\mathrm{pulse}}$ (s)} & \textbf{GATA1\textsubscript{nuc} (mM)} &
\textbf{PU.1\textsubscript{nuc} (mM)} & \textbf{GATA1/PU.1} & \textbf{Fate} &
\textbf{Reversal (s)} \\
\midrule
60  & $\approx$0.000 & 2.563 & 0.00 & Myeloid / reverted & --- \\
90  & 0.002          & 3.616 & 0.00 & Myeloid / reverted & $\sim$1000 \\
\textbf{120} & \textbf{2.170} & \textbf{0.309} & \textbf{7.01} &
\textbf{Erythroid (maintained)} & \textbf{---} \\
\textbf{150} & \textbf{2.704} & \textbf{0.065} & \textbf{41.4} &
\textbf{Erythroid (marginal)} & \textbf{---} \\
300  & 0.019 & 2.761 & 0.01 & Myeloid / reverted & $\sim$2000 \\
500  & $\approx$0.000 & 1.653 & 0.00 & Myeloid / reverted & $\sim$2000 \\
1000 & $\approx$0.000 & 3.769 & 0.00 & Myeloid / reverted & $\sim$2500 \\
1800 & $\approx$0.000 & 0.684 & 0.00 & Myeloid / reverted & $\sim$3000 \\
\bottomrule
\end{tabular}}
\end{table}

The result is non-monotone. A 60~s pulse is insufficient to erase
PU.1 and the erythroid basin is never entered. A 120~s pulse drives
PU.1$_\text{nuc}$ to near zero before withdrawal, leaving GATA1 at approximately
2.0~mM; GATA1 self-activation is sufficient to maintain this state without
ongoing EPOR occupancy. A 150~s pulse also maintains erythroid fate but by a
more marginal trajectory. Pulses of 300~s and longer all drive GATA1 to peak
values above 10~mM---far above the self-maintenance window---and all subsequently
revert; the ghost-attractor reversal timescale scales progressively later with
longer pulses as the GATA1 overshoot grows.

The mechanistic basis is a ghost-attractor effect. Stable post-withdrawal
erythroid fate requires GATA1$_\text{nuc}\in(1.5,\,3)$~mM with
PU.1$_\text{nuc}\lesssim 0.3$~mM at the moment of removal. The 120~s pulse
deposits the system in this window with PU.1 fully suppressed. Longer pulses
drive GATA1 far above this window; the subsequent decay trajectory passes
through the target window at speed while PU.1 simultaneously recovers. The
critical determinant is not whether GATA1 passes through the right concentration
range but whether PU.1 is sufficiently suppressed while it does so.

The lower bound of the stable erythroid window lies between 90 and 120~s; the
upper bound lies between 150 and 300~s. The confirmed erythroid-maintaining
interval is $t_{\mathrm{pulse}} \in [120,\,150]$~s, a double-barrelled
quantitative prediction accessible to continuous live-cell GATA1/PU.1 imaging.

\subsection{The Classical EPO* Threshold Is a Population-Averaging Artefact}
\label{sec:epostar}

The EPO* threshold---the apparent sharp transition in commitment fraction visible
in bulk colony assays---emerges from this model as a consequence of averaging
across stochastic replicates at high copy numbers. When the model is run at
standard (population-scale) initial conditions, the coefficient of variation of
GATA1\textsubscript{nuc} collapses to near zero: all replicates produce nearly
identical trajectories, and a small change in EPO near the threshold switches
the mean trajectory from the myeloid to the erythroid basin, producing the
cliff-like dose-response curve that historically suggested instruction. At
one-hundredth of standard conditions, the same circuit yields 60--70\% erythroid
fraction with enormous replicate variance---the population mean masks independent
draws from a bimodal distribution.

This scale dependence means that the threshold EPO* measured in bulk assays is
the EPO dose that switches the majority of the population mean, which corresponds
to shifting the per-cell commitment probability above 50\%. It is not a cellular
switching threshold: no individual cell has a threshold; each makes a
probabilistic decision influenced by noise internal to the decision layer. The
EPO* is therefore a statistical property of the measurement procedure and
population copy number, not a property of the molecular circuit. This
interpretation is consistent with, and provides a mechanistic explanation for,
the failure of multiple groups to demonstrate all-or-none EPO* switching in
single-cell assays~\cite{hoppe2016,olsson2016}.

\subsection{Validation Against Hoppe \etal\ 2016}
\label{sec:hoppe}

Hoppe \etal~\cite{hoppe2016} used dual-reporter live-cell imaging to track GATA1
and PU.1 protein dynamics in single myeloid progenitors and found that individual
commitment decisions are stochastic with high cell-to-cell variability. To
validate our scale-dependent noise predictions quantitatively, we compare three
outputs of our two simulation regimes against the parameters reported by Hoppe
\etal\ (Fig.~\ref{fig:hoppe_validation}).

\textbf{Commitment fraction.}
At standard initial conditions, every replicate commits to the erythroid
attractor across all EPO concentrations tested (P\,(erythroid)\,=\,100\%,
CV\,$\approx$\,0), consistent with the deterministic-like behaviour seen in
bulk colony assays. At one-hundredth of standard conditions, the same EPO range
yields P\,(erythroid)\,=\,60--70\% (Table~\ref{tab:commitment_D}), statistically
flat across EPO concentrations (Kruskal--Wallis $p>0.05$). This plateau at
sub-maximal commitment matches the wide commitment distributions reported by
Hoppe \etal\ under physiological EPO.

\textbf{Cell-to-cell variability.}
The coefficient of variation of GATA1\textsubscript{nuc} at $t=3600$~s in the
single-cell regime is 0.29, compared to CV\,$\approx$\,0 at population scale.
Hoppe \etal\ reported GATA1 expression CV values of 0.35--0.45 from single-cell
time-lapse data~\cite{hoppe2016}. The simulated CV falls slightly below this
empirical window but is an order of magnitude above the population-scale value,
placing the single-cell regime in the qualitatively correct high-noise
regime.

\textbf{Commitment timing.}
The mean-trajectory GATA1/PU.1 crossing time is near-instantaneous at population
scale, while in the single-cell regime the priming crossover (Section~\ref{sec:priming})
marks a noise-delayed establishment of GATA1 dominance consistent with the
variable commitment latencies observed in live imaging.

\begin{figure}[!htb]
  \centering
  \includegraphics[width=\linewidth]{fig8_hoppe_validation}
  \caption{\textbf{Validation against Hoppe \etal\ 2016.}
  \textbf{A,} Erythroid commitment fraction at $t=3600$~s.
  Population-scale run yields 100\% erythroid at all EPO doses.
  Single-cell run yields 60--70\% (shaded band: Hoppe \etal\ reference
  range~\cite{hoppe2016}).
  \textbf{B,} CV of GATA1\textsubscript{nuc}; population-scale
  CV\,$\approx$\,0; single-cell CV\,=\,0.29 (ochre band: Hoppe 2016 empirical
  range 0.35--0.45).
  \textbf{C,} Mean-trajectory GATA1/PU.1 crossing time at both scales.}
  \label{fig:hoppe_validation}
\end{figure}

\FloatBarrier
\subsection{Nuclear pH Shifts the EPO Commitment Threshold}
\label{sec:ph_sweep}

All prior sweeps held nuclear pH constant at 7.5. To characterise the effect of
nuclear pH on the bifurcation, we ran a factorial EPO--pH sweep: eight EPO
concentrations spanning 0.52--0.65~mM crossed with three nuclear pH values
(7.0, 7.5, 8.0) and GCSF\,=\,1.1~mM, with $N=100$ stochastic replicates per
condition (2400 replicates in total; 24 conditions).

\subsubsection*{pH Shifts the EPO Commitment Threshold}

Erythroid commitment probability increased monotonically with alkaline pH at all
EPO concentrations near the bifurcation. To localise EPO* precisely at each pH,
dedicated fine-grid EPO scans were performed at each of the three pH values
($N=100$--$200$ replicates per condition, $t_{\mathrm{end}}=2000$~s). The EPO
threshold at which P(ERY) crosses 50\% (EPO*) shifted from
0.6125~mM at pH\,=\,7.0 to 0.558~mM at pH\,=\,7.5 and 0.5535~mM at
pH\,=\,8.0 (Table~\ref{tab:ph_threshold}). The
pH\,=\,7.0 to pH\,=\,8.0 shift in EPO* therefore amounts to approximately
0.059~mM. Importantly, a net-flux analysis of the model predicts a single
EPO*\,=\,0.52~mM independent of pH, underestimating the threshold and missing the
pH-dependent shift entirely; the stochastic bistability of the switch is invisible
to mean-field analysis.

\begin{table}[ht]
\caption{EPO commitment threshold (EPO*) by nuclear pH, determined from
dedicated fine-grid EPO scans ($N=100$--$200$ per condition, GCSF\,=\,1.1~mM,
$t_{\mathrm{end}}=2000$~s). Mean-field (net-flux) prediction shown for comparison.}
\label{tab:ph_threshold}
\centering
\small
\begin{tabularx}{\linewidth}{@{}Xcc@{}}
\toprule
\textbf{Nuclear pH} & \textbf{EPO* (stochastic, mM)} & \textbf{EPO* (net-flux, mM)} \\
\midrule
7.0 & 0.6125 & 0.52 \\
7.5 & 0.558  & 0.52 \\
8.0 & 0.5535 & 0.52 \\
\bottomrule
\end{tabularx}
\end{table}

The pH mechanism acts through the inhibition half-maximum constant governing
mutual repression between GATA1 and PU.1 transcription:
\begin{equation}
K_{\mathrm{inh}}(\mathrm{pH}) = 8.0 \times 10^{\,0.5(\mathrm{pH} - 7.5)}
\ \mathrm{mM}
\label{eq:pHhill}
\end{equation}
At pH\,=\,7.0, $K_{\mathrm{inh}} = 4.5$~mM; at pH\,=\,7.5 (baseline),
$K_{\mathrm{inh}} = 8.0$~mM; at pH\,=\,8.0, $K_{\mathrm{inh}} = 14.2$~mM.
A lower $K_{\mathrm{inh}}$ strengthens cross-inhibition in both directions:
near the bifurcation where PU.1 dominates, acidic pH strengthens PU.1-mediated
suppression of GATA1, raising the EPO barrier to commitment---consistent with
the higher EPO* at pH\,=\,7.0.

\subsubsection*{Fate Reversal at a Single EPO Concentration}

The practical consequence of the pH-dependent EPO* shift is dramatic near the
bifurcation. At EPO\,=\,0.57~mM---a single concentration---the erythroid
fraction is 4\% at pH\,=\,7.0 but 95\% at pH\,=\,8.0: the same cytokine dose
produces opposite majority fates at the two pH extremes. This finding constitutes
an experimentally testable prediction: adjusting nuclear pH while holding EPO
fixed near the critical value should switch the population from predominantly
myeloid to predominantly erythroid.

\subsubsection*{More Than Half of Conditions Are Genuinely Bistable}

To determine whether the bimodal fate distributions reflect genuine stochastic
bistability rather than noise around a unimodal attractor, we applied the
bimodality coefficient (BC) to the GATA1\textsubscript{nuc}/PU.1\textsubscript{nuc}
ratio distribution at the end of each simulation. A value BC\,$>$\,0.555 indicates
departure from unimodality; 13 of the 24 conditions exceeded this threshold, with
the highest BC values (BC\,$>$\,0.88) clustered near the bifurcation at each pH.
At EPO\,=\,0.57~mM---below EPO* for pH\,=\,7.0 (0.6125~mM) but above EPO*
for pH\,=\,7.5 (0.558~mM) and pH\,=\,8.0 (0.5535~mM)---all three pH
conditions were bistable. Conditions far above or below EPO* were
unimodal, with all replicates collapsing to the same attractor.

\begin{table}[ht]
\caption{Bistability summary from the EPO--pH factorial sweep ($N=100$ per
condition, GCSF\,=\,1.1~mM). Bimodality coefficient (BC) $>$\,0.555 indicates
stochastic bistability. BC\,$>$\,0.88 marks conditions with the sharpest
bimodal separation. 13 of 24 conditions are bistable; all 24 conditions
tabulated in supplementary material.}
\label{tab:bistability}
\centering
\small
\begin{tabularx}{\linewidth}{@{}cccc@{}}
\toprule
\textbf{EPO (mM)} & \textbf{pH} & \textbf{P(ERY)} & \textbf{BC} \\
\midrule
0.52 & 7.0  & 0.06 & 0.29 \\
0.52 & 7.5  & 0.08 & 0.34 \\
0.52 & 8.0  & 0.21 & 0.72 \\
0.57 & 7.0  & 0.04 & 0.89 \\
0.57 & 7.5  & 0.48 & 0.95 \\
0.57 & 8.0  & 0.95 & 0.91 \\
0.61 & 7.0  & 0.42 & 0.93 \\
0.61 & 7.5  & 0.78 & 0.88 \\
0.61 & 8.0  & 0.99 & 0.31 \\
0.65 & 7.0  & 0.91 & 0.47 \\
0.65 & 7.5  & 0.97 & 0.29 \\
0.65 & 8.0  & 1.00 & 0.18 \\
\bottomrule
\end{tabularx}
\end{table}

\subsubsection*{Attractor Depth Is pH-Dependent in the Opposite Direction}

A complementary experiment ran EPO\,=\,1.0~mM with GCSF\,$\approx$\,0 at
each of the three pH values ($N=50$ replicates per pH, 150 replicates total),
placing all cells deep in the erythroid attractor (P\,(ERY)\,=\,100\% at all
three pH values). Residual nuclear PU.1 in committed cells increased
monotonically with alkaline pH: mean PU.1$_\text{nuc}$ was 0.000092~mM at
pH\,=\,7.0, 0.004~mM at pH\,=\,7.5, and 0.006~mM at pH\,=\,8.0
(Table~\ref{tab:attractor_depth}). Acidic conditions produced the deepest
erythroid attractor---once committed, GATA1 more completely suppresses PU.1 at
low pH, because the smaller $K_{\mathrm{inh}}$ makes cross-inhibition more
potent in the GATA1-dominant state.

\begin{table}[ht]
\caption{Residual nuclear PU.1 in committed cells at EPO\,=\,1.0~mM,
GCSF\,$\approx$\,0 ($N=50$ per pH). Smaller PU.1$_\text{nuc}$ indicates a
deeper erythroid attractor. All residual values are below 0.01~mM against a
GATA1$_\text{nuc}$ background of approximately 9~mM---the pH effect on attractor
depth is a second-order modulation.}
\label{tab:attractor_depth}
\centering
\small
\begin{tabularx}{\linewidth}{@{}Xccc@{}}
\toprule
\textbf{Nuclear pH} & \textbf{PU.1$_\text{nuc}$ mean (mM)} & \textbf{$\pm$\,SD}
& \textbf{Max (mM)} \\
\midrule
7.0 & 0.000092 & 0.000649 & 0.0046 \\
7.5 & 0.003873 & 0.021241 & 0.143  \\
8.0 & 0.006114 & 0.030238 & 0.156  \\
\bottomrule
\end{tabularx}
\end{table}

The two pH experiments together reveal two separable regimes governed by the
same parameter: at near-bifurcation EPO, alkaline pH lowers EPO* and increases
commitment probability; in cells already committed at high EPO, acidic pH
produces marginally more complete PU.1 silencing. Nuclear pH therefore acts as a
weak, bidirectional modulator of erythroid commitment---regime-dependent in
direction---rather than a unidirectional rheostat, and its effect on the EPO
threshold is quantitatively distinct from what mean-field analysis predicts.
